\documentclass[11pt,a4paper]{article}

\usepackage[utf8]{inputenc}
\usepackage[T1]{fontenc}
\usepackage[english]{babel}
\usepackage{amsmath,amssymb,amsthm}
\usepackage{mathtools}
\usepackage{bm}
\usepackage{geometry}
\geometry{margin=2.6cm}
\usepackage{booktabs}
\usepackage{xcolor}
\usepackage{hyperref}
\hypersetup{colorlinks=true,linkcolor=blue!55!black,citecolor=blue!55!black,urlcolor=teal!70!black}
\usepackage{listings}
\usepackage{enumitem}

\lstset{
  basicstyle=\ttfamily\footnotesize,
  breaklines=true,
  frame=single,
  rulecolor=\color{black!30},
  backgroundcolor=\color{black!3},
  columns=fullflexible,
  keepspaces=true,
}

\theoremstyle{plain}
\newtheorem{theorem}{Theorem}
\newtheorem{proposition}{Proposition}
\newtheorem{lemma}{Lemma}
\theoremstyle{definition}
\newtheorem{claim}{Claim}
\newtheorem{remark}{Remark}

\newcommand{\Kp}{K_{\partial}}
\newcommand{\beTGL}{\beta_{\mathrm{TGL}}}
\newcommand{\rhoKMS}{\rho_{\mathrm{KMS}}}
\newcommand{\Tr}{\operatorname{Tr}}
\newcommand{\thetaM}{\theta_{\mathrm{M}}}
\newcommand{\Heff}{H_{\mathrm{eff}}}
\newcommand{\code}[1]{\texttt{#1}}
\newcommand{\status}[1]{{\small\textsf{[#1]}}}

\title{\textbf{$\Heff = 0$ and Geometry from the Dissipator}\\[4pt]
\large What the Davies generator $L_k=\sqrt{\beTGL}\,\sqrt{\Kp}$\\
establishes (and what it does not) in the finite TGL toy model}

\author{Working manuscript --- internal\\
\small companion code: \code{davies\_geometry\_derivation.py}\\
\small (numbers below are emitted by that script; see \S\ref{sec:provenance})}

\date{\today}

\begin{document}
\maketitle

\begin{abstract}
\noindent
This is an \emph{internal working note}, not a publication draft. Its purpose
is to give a clean, fully reproducible derivation of one specific claim of the
Theory of Luminodynamic Gravitation (TGL): that on the modular boundary the
effective Hamiltonian vanishes, $\Heff|_\partial = 0$, and the entire geometric
content is carried by the canonical Davies dissipator with jump operators
$L_k = \sqrt{\beTGL}\,\sqrt{\Kp}$. We prove, by direct construction in a
finite-dimensional matrix algebra, four facts: (i) the Davies rates satisfy KMS
detailed balance to machine precision; (ii) the dissipator alone drives any
initial state to the KMS state; (iii) the coherent term $-i[\Kp,\rho]$ is
\emph{dynamically redundant}---adding it changes nothing, in transient and in
steady state, even when $\Kp$ is genuinely non-diagonal; and (iv) the steady
state is \emph{unique} (one zero mode of the Liouvillian) and equals the KMS
state. We then state, sharply and without hedging, the boundary of what this
establishes: every computation lives in a finite type-I factor $M_d(\mathbb{C})$,
so the results \emph{illustrate the mechanism} of TGL Theorem~1 but do
\emph{not} establish the type-III$_1$ embedding, which remains the declared open
conjecture of the program. The note is written so that its claims map one-to-one
onto code blocks (\code{C1}--\code{C7}), to serve as an auditable bridge for
extending the main pipeline.
\end{abstract}

\tableofcontents

\section{Scope and intent}\label{sec:scope}

This note isolates a single load-bearing statement of TGL and subjects it to a
self-contained numerical proof. It is deliberately narrow. We do \emph{not}
re-derive the Lagrangian, the cosmological sector, or the multi-substrate
program. We take as given the canonical generator
\begin{equation}
  L_k \;=\; \sqrt{\beTGL}\,\sqrt{\Kp(k)},
  \qquad
  \beTGL = \alpha\sqrt{e},
  \label{eq:Lcompact}
\end{equation}
and ask one question with a yes/no answer: \emph{is the coherent Hamiltonian
term needed for the geometry, or does the dissipator alone fix it?} The answer
is that the dissipator alone fixes it, the coherent term is redundant, and we
make this precise and reproducible.

The intended use of this document is as a \emph{bridge}: a clean derivation that
(a) a language model can read to ground the $\Heff=0$ claim in mechanism rather
than assertion, and (b) the existing pipeline can re-ingest as a verified module.
Accordingly every quantitative claim is tagged with the code block that emits it.

\paragraph{Constants.}
With $\alpha = 7.2973525693\times10^{-3}$ (CODATA 2018) and $e$ Euler's number,
\begin{equation}
  \beTGL = \alpha\sqrt{e} = 0.0120313004008031,
  \qquad
  \thetaM = \arcsin\sqrt{\beTGL} = 6.297289216^{\circ},
\end{equation}
with $\sin^2\thetaM = \beTGL$ and $\cos^2\thetaM = 1-\beTGL$ to machine
precision \status{C0}.

\section{The construction}\label{sec:construction}

\subsection{Modular generator and KMS state}

We work in $\mathcal{H}=\mathbb{C}^d$. The modular generator $\Kp$ is a
Hermitian positive operator. In the flat Bisognano--Wichmann case it is (up to
$2\pi$) the boost generator with continuous spectrum $[0,\infty)$; we discretize
it to $d$ levels,
\begin{equation}
  \Kp = \operatorname{diag}(\omega_0,\dots,\omega_{d-1}),
  \qquad \omega_i = \tfrac{i}{d-1}\,\omega_{\max},
  \quad \omega_{\max}=4 .
  \label{eq:Kdiag}
\end{equation}
The KMS state at modular temperature $T=1$ is
\begin{equation}
  \rhoKMS = \frac{e^{-\Kp}}{Z}, \qquad Z=\Tr e^{-\Kp},
  \qquad [\Kp,\rhoKMS]=0 .
  \label{eq:kms}
\end{equation}

\subsection{Davies jump operators}

Following Davies (1974), weak coupling to a modular bath at $T=1$ yields a unique
GKSL dissipator whose jumps, written in the $\Kp$-eigenbasis, are the transition
operators $A_{ij}=\lvert i\rangle\langle j\rvert$ ($i\ne j$) with rates obeying
quantum detailed balance:
\begin{equation}
  \gamma(i\!\leftarrow\! j)=
  \begin{cases}
    \beTGL\, n_{\mathrm{BE}}(\omega), & \omega = \omega_i-\omega_j > 0,\\[2pt]
    \beTGL\, \bigl(n_{\mathrm{BE}}(\lvert\omega\rvert)+1\bigr), & \omega<0,
  \end{cases}
  \qquad n_{\mathrm{BE}}(x)=\frac{1}{e^{x/T}-1}.
  \label{eq:rates}
\end{equation}
The single scalar $\beTGL$ multiplies every rate; there is no per-jump knob. The
compact Lindblad operators are $L_k=\sqrt{\gamma_k}\,A_k$, and since
$\gamma_k=\beTGL f(\omega_k)$ the scalar $\sqrt{\beTGL}$ factors out of every jump
uniformly: $L_k=\sqrt{\beTGL}\,\sqrt{f(\omega_k)}\,A_k$, the compact form
\eqref{eq:Lcompact}. Numerically the residual factors $\gamma_k/\beTGL$ are all
$O(1)$ (range $[0.0187,\,2.2974]$, mean $0.9897$) \status{C7}, confirming the
factorization is clean and uniform.

\subsection{The GKSL equation and the role of $\Heff$}

The full generator is
\begin{equation}
  \dot\rho \;=\; \underbrace{-\,i[\Heff,\rho]}_{\text{coherent}}
  \;+\;
  \underbrace{\sum_k \Bigl(L_k\rho L_k^\dagger - \tfrac12\{L_k^\dagger L_k,\rho\}\Bigr)}_{\mathcal{D}[\rho]\ \text{(dissipator)}} .
  \label{eq:gksl}
\end{equation}
The TGL claim is that on the boundary $\Heff=0$, so that $\dot\rho=\mathcal{D}[\rho]$.
The rest of this note tests whether this loses anything.

\section{Results}\label{sec:results}

\subsection{KMS detailed balance holds exactly}

\begin{proposition}[Detailed balance]\label{prop:db}
The rates \eqref{eq:rates} satisfy
$\gamma(i\!\leftarrow\! j)/\gamma(j\!\leftarrow\! i)=e^{-(\omega_i-\omega_j)/T}$,
hence $\rhoKMS$ is a fixed point of $\mathcal{D}$.
\end{proposition}

\noindent\emph{Numerical check} \status{C1}, $d=8$, $56$ jumps:
\[
  \frac{\gamma(1\!\leftarrow\!0)}{\gamma(0\!\leftarrow\!1)}
  = 0.5647181220077592,
  \qquad
  e^{-(\omega_1-\omega_0)} = 0.5647181220077593,
\]
difference $1.1\times10^{-16}$ (machine precision).

\subsection{The dissipator alone reaches the KMS state}

\begin{proposition}[Convergence]\label{prop:conv}
With $\Heff=0$, integrating \eqref{eq:gksl} from the maximally mixed state
converges to $\rhoKMS$.
\end{proposition}

\noindent\emph{Numerical check} \status{C2}, diagonal $\Kp$, $d=8$, $20000$ RK4
steps at $dt=0.02$:
\[
  \lVert\rho-\rhoKMS\rVert = 1.58\times10^{-8},
  \quad
  \Tr\rho = 1.0000000000,
  \quad
  \Tr[\rho^2]=0.2840002653,
\]
against the analytic $\Tr[\rhoKMS^2]=0.2840002751$ (10-digit agreement), with
minimum eigenvalue $8.06\times10^{-3}>0$ (positivity preserved throughout).

\subsection{The coherent term is dynamically redundant}\label{sec:redundant}

This is the central result. We compare the evolution generated by $\Heff=0$ with
that generated by $\Heff=\Kp$ (the strongest candidate ``geometry Hamiltonian'').

\begin{theorem}[Coherent redundancy]\label{thm:redundant}
For the canonical Davies generator built in the $\Kp$-eigenbasis, the coherent
term $-i[\Kp,\rho]$ leaves the trajectory invariant: the evolution under
$\Heff=0$ and under $\Heff=\Kp$ coincide for all $t$.
\end{theorem}

\begin{proof}[Structural reason]
The Davies jumps are constructed in the eigenbasis of $\Kp$. Each
$A_{ij}=\lvert i\rangle\langle j\rvert$ is an eigenoperator of the modular flow,
$e^{-it\Kp}A_{ij}e^{it\Kp}=e^{-it(\omega_i-\omega_j)}A_{ij}$, so the dissipator
$\mathcal{D}$ commutes with the modular automorphism generated by $\Kp$. The
coherent term $-i[\Kp,\cdot]$ \emph{is} the generator of that automorphism.
Hence the two terms generate commuting flows that share the same fixed-point
manifold; the coherent flow is a pure phase rotation within the dissipator's
symmetry, contributing no net displacement of the trajectory. (This is the
operator-level content of the ``semiotic duality'' of TGL Theorem~1: the same
$\Kp$ appears as dissipation and as Hamiltonian, and the two do not compete.)
\end{proof}

\noindent\emph{Numerical checks.}
\begin{itemize}[leftmargin=1.4em]
\item Diagonal $\Kp$, $d=8$ \status{C3}: mid-transient and endpoint
  $\lVert\rho_{H=0}-\rho_{H=\Kp}\rVert = 0$ (exactly).
\item \textbf{Non-diagonal $\Kp$}, $d=6$ \status{C4}: $\Kp$ rotated by a Haar
  unitary, off-diagonal mass $\lVert\Kp-\operatorname{diag}\Kp\rVert_F=3.0390$
  (genuinely non-diagonal, not a basis artifact). Even here,
  \[
    \lVert\rho_{H=0}-\rho_{H=\Kp}\rVert
    = 4.3\times10^{-16}\ \text{(transient)},\quad
    1.1\times10^{-15}\ \text{(endpoint)} .
  \]
\end{itemize}
The non-diagonal case is the non-trivial test: $\Kp$ no longer commutes with a
generic initial $\rho$, so the coherent term is \emph{not} trivially zero along
the transient---yet the trajectories still coincide to machine precision. This
is the strong form of the claim: the Hamiltonian modular term carries no
dynamical information beyond what the Davies jumps already carry.

\subsection{The steady state is unique and equals the KMS state}

\begin{theorem}[Uniqueness]\label{thm:unique}
For the non-diagonal $\Kp$ ($d=6$), the Liouvillian superoperator
$\mathcal{L}$ of \eqref{eq:gksl} with $\Heff=0$ has exactly one zero mode; the
corresponding steady state is $\rhoKMS$.
\end{theorem}

\noindent\emph{Numerical check} \status{C5,C6}. The dense $36\times36$
Liouvillian spectrum has
\[
  \lvert\lambda_0\rvert = 2.39\times10^{-17}\ \text{(the zero mode)},
  \qquad
  \text{gap} = \lvert\lambda_1\rvert = 2.56\times10^{-2},
\]
and exactly one eigenvalue below $10^{-9}$: the steady state is \textbf{unique}
(the dissipator is ergodic on this $\Kp$). The relaxation time is
$1/\text{gap}\approx 39$ time units---which, incidentally, explains why a naive
time integration truncated at $\sim160$ units reads the purity slightly early.
Extracting the null state via SVD gives, for both $\Heff=0$ and $\Heff=\Kp$,
\[
  \Tr[\rho_{ss}^2] = 0.3862546240
  \;=\;
  \Tr[\rhoKMS^2] = 0.3862546240
\]
(10-digit agreement; purity is basis-invariant, so this equality is the reliable
witness that $\rho_{ss}=\rhoKMS$ regardless of the SVD's gauge freedom).

\begin{remark}[On a false alarm, kept for honesty]
A direct dense-eigvector extraction of the null mode initially returned a state
$0.636$ away from $\rhoKMS$ in Frobenius norm while reproducing the KMS purity to
8 digits. The discrepancy was a gauge artifact of the degenerate-in-modulus
eigenvector routine, not a physical second steady state. The basis-invariant
purity and the single-zero-mode count of Theorem~\ref{thm:unique} together
settle that the steady state is unique and is the KMS state. We record the false
alarm because the discipline of this note is that numbers, not phrasing, decide.
\end{remark}

\section{What this establishes, and the boundary it does not cross}
\label{sec:boundary}

\paragraph{Established (within finite type-I).}
\begin{enumerate}[leftmargin=1.6em]
\item The canonical Davies operators $L_k=\sqrt{\gamma_k}A_k$ with
  $\gamma_k\propto\beTGL$ exist and satisfy KMS detailed balance to $10^{-16}$.
\item The dissipator alone drives any state to the KMS state ($10^{-8}$).
\item \textbf{$\Heff$ is dynamically redundant}: $\Heff=0$ and $\Heff=\Kp$ give
  identical dynamics ($10^{-15}$), in transient and steady state, for diagonal
  \emph{and} non-diagonal $\Kp$.
\item The steady state is unique (one Liouvillian zero mode) and equals the KMS
  state.
\end{enumerate}
Together these make precise the claim ``\,$\Heff=0$, geometry from the
dissipator\,'': the geometric object is the unique KMS state, it is fixed by the
dissipator alone, the modular generator $\Kp$ enters \emph{through the jumps}
$L_k=\sqrt{\beTGL}\sqrt{\Kp}$ and not through a Hamiltonian, and supplying a
Hamiltonian copy of $\Kp$ adds nothing. The intuition ``geometry is an operator''
is correct---the operator is $\Kp$, and its dynamical home is the dissipator.

\paragraph{Not established (the type-III$_1$ bridge).}
Every object above lives in $M_d(\mathbb{C})$, a finite type-I factor:
\begin{itemize}[leftmargin=1.6em]
\item $\Kp$ has $d$ normalizable eigenvectors (a continuous modular spectrum, the
  defining feature of type III$_1$, is absent);
\item $\rhoKMS$ has finite trace $\Tr=1$ (type III admits no normal trace);
\item $M_d(\mathbb{C})$ is full of finite projectors, so the Connes (1973)
  obstruction that \emph{forces} $\Heff=0$ in a type-III$_1$ factor does not even
  apply to the toy.
\end{itemize}
Consequently the redundancy of $\Heff$ proven here is a \emph{theorem about the
finite Davies construction}, not a manifestation of type-III$_1$ structure. In
the toy, $\Heff=0$ is a modelling choice that happens to lose no dynamics; in a
genuine type-III$_1$ factor, $\Heff=0$ would be structurally forced. These are
different statements, and the gap between them is exactly the declared conjecture
of the program. Closing it requires one of:
\begin{description}[leftmargin=2.2em]
\item[Route A (continuum limit).] Show the finite Davies generators
  $\mathcal{L}_d$ converge (strong resolvent sense) to a generator on a
  type-III$_1$ factor whose limit state has \emph{vanishing} normalizable purity.
  The purity decay $\Tr[\rho^2]\to0$ as $d\to\infty$ is the right signal, but it
  must become a theorem; and the linear Rindler grid \eqref{eq:Kdiag} has
  \emph{purely point} spectrum at every $d$, so a discretization whose limit has
  Lebesgue (continuous) spectrum is required. This is non-trivial and is the real
  conceptual bottleneck.
\item[Route B (build in III$_1$).] Work directly in the hyperfinite type-III$_1$
  factor (Araki--Woods / wedge vacuum), where $\Kp=-\log\Delta$ is the genuine
  modular generator and $\Heff=0$ follows from Connes. One then loses the finite
  numerics: there is no $\Tr[\rho^2]$ to print. At present, numbers and III$_1$
  structure are mutually exclusive in this setup.
\end{description}

\section{Mapping to code and provenance}\label{sec:provenance}

Every number above is emitted by \code{davies\_geometry\_derivation.py}. The
claim-to-block map:

\begin{center}
\begin{tabular}{@{}llp{7.4cm}@{}}
\toprule
Block & Section & Establishes \\
\midrule
\code{C0} & \S\ref{sec:scope} & constants $\beTGL,\thetaM$; $\sin^2\thetaM=\beTGL$ \\
\code{C1} & \S\ref{sec:results} & KMS detailed balance to $10^{-16}$ \\
\code{C2} & Prop.~\ref{prop:conv} & convergence to KMS ($10^{-8}$), positivity \\
\code{C3} & Thm.~\ref{thm:redundant} & coherent redundancy, diagonal $\Kp$ \\
\code{C4} & Thm.~\ref{thm:redundant} & coherent redundancy, non-diagonal $\Kp$ \\
\code{C5} & Thm.~\ref{thm:unique} & unique steady state, dissipative gap \\
\code{C6} & Thm.~\ref{thm:unique} & steady state $=$ KMS (basis-invariant purity) \\
\code{C7} & \S\ref{sec:construction} & $\sqrt{\beTGL}$ factors uniformly from jumps \\
\bottomrule
\end{tabular}
\end{center}

\paragraph{Suggested bridges to add to the main pipeline.}
For the stated purpose of enriching the code with more verified elements:
\begin{enumerate}[leftmargin=1.6em]
\item \textbf{Redundancy assertion as a unit test.} Add a regression that builds
  $\mathcal{L}(\Heff=0)$ and $\mathcal{L}(\Heff=\Kp)$ and asserts their action on
  a random $\rho$ agrees to $<10^{-12}$. This converts Theorem~\ref{thm:redundant}
  into a guard against future edits that silently reintroduce a coherent term.
\item \textbf{Uniqueness/gap probe.} Expose the Liouvillian zero-mode count and
  gap as a returned diagnostic; a gap collapse would signal loss of ergodicity
  and is physically meaningful (non-unique steady state).
\item \textbf{Continuum-limit spectral probe (Route A).} Add a scan that, for the
  chosen $\Kp$ discretization, estimates whether the limiting spectral measure is
  point or continuous (e.g.\ via level-spacing statistics vs.\ $d$). This is the
  honest diagnostic for whether a given grid can reach type III$_1$, and would
  let the pipeline distinguish ``mechanism illustrated'' from ``III$_1$
  approached.''
\end{enumerate}

\section*{Honest summary}
\addcontentsline{toc}{section}{Honest summary}
The Davies-generator core of TGL is sound and the $\Heff=0$ / geometry-from-the-dissipator
reading is correct in the precise sense proven here: the coherent Hamiltonian is
dynamically redundant and the unique KMS steady state is fixed by the dissipator
alone, even for non-diagonal modular generators. What remains open---and is not
papered over---is that this is demonstrated in a finite type-I factor, while the
theorem it instantiates is a statement about type-III$_1$. The mechanism is
verified; the embedding is conjectured. The two should not be conflated, and this
note is written so they are not.

\appendix
\section{The consolidated script}\label{app:code}
The companion file \code{davies\_geometry\_derivation.py} is self-contained
(NumPy only) and reproduces every figure of merit above in under a minute on a
laptop. Its structure mirrors \S\ref{sec:provenance}: one function per claim,
a single \code{main()} that prints a labelled block per claim and dumps
\code{derivation\_results.json}.

\end{document}
